% Compile: lualatex

\documentclass[fontsize=14pt]{scrartcl}



\usepackage[preset=modeling]{bmstu-report}

% \usepackage{fontspec}
% % Последний язык в списке — основной (russian); english — для англ. заголовков.
% \usepackage[english,russian]{babel}
% \usepackage{amsmath}
% \defaultfontfeatures{Ligatures = TeX}
% \setmainfont{CMU Serif}[
%   Extension      = .otf,
%   UprightFont    = cmunrm,
%   BoldFont       = cmunbx,
%   ItalicFont     = cmunti,
%   BoldItalicFont = cmunbi,
% ]
% \setsansfont{CMU Sans Serif}[
%   Extension      = .otf,
%   UprightFont    = cmunss,
%   BoldFont       = cmunsx,
%   ItalicFont     = cmunsl,
%   BoldItalicFont = cmunbl,
% ]
% \setmonofont{CMU Typewriter Text}[
%   Extension   = .otf,
%   UprightFont = cmuntt,
%   BoldFont    = cmuntb,
%   ItalicFont  = cmunit,
% ]

\usepackage{graphicx}
\usepackage{tikz}
\usetikzlibrary{positioning,calc}
\usetikzlibrary{backgrounds}
\pgfdeclarelayer{background}
\pgfsetlayers{background,main}

\usepackage{amssymb}
\newcommand{\bigsqcap}{\mathop{\raisebox{-0.15ex}{\scalebox{1.6}{$\sqcap$}}}\displaylimits}

\newcommand{\Lang}{\mathcal{L}}
\DeclareMathOperator{\Fix}{Fix}
\DeclareMathOperator{\StdMor}{StdMor}
\newcommand{\SP}{\mathrm{SP}}
\newcommand{\SOB}{\mathrm{SOB}}
\newcommand{\Cons}{\mathsf{Cons}}
\newcommand{\bigboxtimes}{\mathop{\text{\raisebox{-0.2ex}{\scalebox{1.6}{$\boxtimes$}}}}}

% --- символы редуцированного и декартова произведения, как в StringIntervalsArxiv.tex ---
\newcommand{\Cartesian}[1]{\mathbin{\begin{tikzpicture}[x=#1ex,y=#1ex,line width=(#1)/4.0*1ex] \draw (-0.7,-0.9) -- (0.7,0.5) (-0.7,0.5) -- (0.7,-0.9); \draw (0,-0.2) circle (1.41);\end{tikzpicture}}}
\newcommand{\tcartesian}{\raisebox{-0.2ex}{\,$\Cartesian{0.6}$\,}}
\newcommand{\Reduced}[1]{\mathbin{\begin{tikzpicture}[x=#1ex,y=#1ex,line width=(#1)/4.0*1ex] \draw (-0.7,-0.7) -- (0,0.7) (0.7,-0.7) -- (0,0.7); \draw (0,0) circle (1.41);\end{tikzpicture}}}
\newcommand{\treduced}{\raisebox{-0.2ex}{\,$\Reduced{0.6}$\,}}
\newcommand{\bigreduced}[2]{\mathop{\raisebox{-0.8ex}{$\Reduced{1.2}$}}_{#1}^{#2}}

\newlength{\tabparindent}
% Текст в ячейке таблицы с обычными абзацами и красной строкой (для \texttt{tabularx}).
\newcommand{\tabparblock}[1]{%
  \begin{minipage}[t]{\linewidth}%
    \setlength{\parindent}{\tabparindent}%
    \setlength{\parskip}{0pt plus .2pt}%
    \strut\ignorespaces
    #1%
  \end{minipage}%
}


\setminted{
  breakanywhere,
  breakautoindent,
}

\providecommand{\prog}[1]{{\ttfamily\small#1}}
\providecommand{\anonsection}[1]{%
  \cleardoublepage
  \phantomsection
  \addcontentsline{toc}{section}{\protect\numberline{}#1}%
  \section*{#1}\vspace*{2.5ex}%
}
\providecommand{\sectionbreak}{\clearpage}

\renewcommand{\cftsecleader}{\cftdotfill{\cftdotsep}}
\renewcommand{\cftsecfont}{\normalfont\large}
\renewcommand{\cftsubsecfont}{\normalfont\large}
\setlength{\cftsecindent}{0pt}
\setlength{\cftsubsecindent}{0pt}
\setlength{\cftsubsubsecindent}{0pt}

\makeatletter
\@ifundefined{abs}{%
  \DeclarePairedDelimiter\abs{\lvert}{\rvert}%
}{}
\@ifundefined{norm}{%
  \DeclarePairedDelimiter\norm{\lVert}{\rVert}%
}{}
\providecommand{\textvisiblespace}{%
  \mbox{\kern.06em\vrule \@height.3ex}%
  \vbox{\hrule \@width.3em}%
  \hbox{\vrule \@height.3ex}}%
\newsavebox{\spacebox}
\begin{lrbox}{\spacebox}
\verb*! !
\end{lrbox}
\providecommand{\aspace}{\usebox{\spacebox}}
\makeatother


\usepackage[backend=biber,
    bibstyle=gost-numeric,
    citestyle=nature]{biblatex}
\addbibresource{biblio.bib}

% Список источников (не подстрочные/затекстовые ссылки): ГОСТ Р 7.0.100—2018, biblatex-gost §5.6
\AtBeginBibliography{%
  \makeatletter
  \toggletrue{bbx:gostbibliography}%
  \renewcommand*{\mkgostheading}[1]{#1}%
  % Заглавие, затем « / » и все авторы (прямой шрифт), без курсива в «заголовке»
  \renewbibmacro*{author}{}%
  \renewbibmacro*{byauthor}{%
    \ifboolexpr{
      test \ifuseauthor
      and
      not test {\ifnameundef{author}}
    }
      {\setrespdelim%
       \printnames[byauthor]{author}}
      {}}%
  \makeatother
}

\begin{document}

\setcounter{page}{2}

\renewcommand\contentsname{\hfill{\normalfont{СОДЕРЖАНИЕ}}\hfill}
\tableofcontents

% \section*{ВВЕДЕНИЕ}
\anonsection{СПИСОК ОБОЗНАЧЕНИЙ И СОКРАЩЕНИЙ}

В работе используются следующие обозначения и сокращения: 

\vspace{0.3cm}

\noindent
\begin{tabular}{@{}>{\setlength{\parindent}{0pt}\raggedright\arraybackslash}p{3.3cm}@{\quad---\quad}>{\setlength{\parindent}{0pt}\raggedright\arraybackslash}p{\dimexpr\linewidth-2.2cm-4em\relax}@{}}
  $\varepsilon$ & пустое слово \\
  $\Sigma$      & конечный алфавит \\
  $\Sigma^*$    & множество всех строк над $\Sigma$ \\
  $\Sigma^+$    & множество всех непустых строк над $\Sigma$ \\
  $\sigma_1 \circ \sigma_2$ & композиция морфизмов $\sigma_1(\sigma_2(x))$ \\
  $\sigma^{-1}(x)$ & прообраз слова $x$ в $\sigma$ \\
  $\sigma|_{A}$ & ограничение морфизма $\sigma$ на множество $A$ \\
  $im(\sigma)$ & образ морфизма $\sigma$ \\
  $\Lang(S)$    & язык, задаваемый $S$ \\
  LTL & локально тестируемый язык (locally testable language) \\
  LTL-свойство & свойство, определяющее локально тестируемый язык  \\
  $\SP(\sigma)$ & класс строковых свойств, отвечающих стандартному морфизму $\sigma$ \\
  ДКА & детерминированный конечный автомат
\end{tabular}

\newpage

\anonsection{ВВЕДЕНИЕ}

В данной работе рассматривается и дополняется теоретическая составляющая 
подхода к абстрактной интерпретации строковых переменных, основанного на 
использовании систем линейных уравнений в словах определенного вида. 

Представленный в статье~\cite{NepeivodoaAfanasyev2025} подход к абстрактной интерпретации строк основан на идее
отслеживания префиксов, суффиксов и подслов (инфиксов), что позволяет довольно
эффективно проводить требуемые операциями над абстрактными строковыми объектами. С целью 
повышения точности анализа предлагается также отслеживать образы строк при некоторых строковых морфизмах.
Этот подход позволит отслеживать непрерывные подстроки определенного вида: например,
морфизм, отображающий символы-цифры в один символ, позволит сохранить информацию о непрерывных
подстроках-числах, в то время как без отслеживания образов будет иметься информация только о разрозненном вхождении 
некоторых цифр в строку.

Абстракция образа строки при некотором морфизме представляется понятием строкового свойства. Именно 
набор строковых свойств образует строковый объект, выступающий абстракцией над строковой переменной.
Однако одно и то же множество конретизации могут описывать несколько строковых объектов, неравных между собой
при явном сопоставлении компонент. Поэтому возникает задача нормализации строкового объекта в целом,
что в том числе позволит более эффективно проводить статический анализ, поскольку некоторые гарантированные
компоненты будут добавлены в рассмотрение на этапе нормализации.

Сопутствующей веткой исследование является рассмотрение морфизмов, которые целесообразно использовать
в инструменте анализа в целом, поскольку они напрямую определяют взаимодействие строковых свойств во внутренней структуре строкового объекта.
Предлагается рассматривать класс так называемых стандартных морфизмов --- идемпотентных 
регрессивных алфавитных морфизмов.

Таким образом, целью работы является исследование нормализации строковых объектов,
которое включает задачи:

\begin{enumerate}
\item Формализовать нормализацию строкового свойства.
\item Исследовать и описать решетки строковых свойств и решетки стандартных морфизмов.
\item Сформулировать концепции нормализации строкового объекта в целом.
\end{enumerate}



\section{Строковые свойства и строковые объекты}
\label{sec:string-properties}

Уравнением в словах называется равенство вида $\Phi = \Psi$, где $\Phi$ и $\Psi$ --- слова над объединенным 
алфавитом $\Sigma \cup X$, где $X$ --- множество переменных.

В данной работе будем рассматривать системы линейных уравнений в словах вида
\begin{equation}
  \label{eq:system-of-equations}
\begin{cases}
    Z = c \\
    Z = p W_0\\
    Z = W_1 s\\
    Z = X_1 w_1 Y_1 \\
    Z = X_2 w_2 Y_2 \\
    \vdots \\
    Z = X_n w_n Y_n \\
\end{cases},
\end{equation}

где, возможно, отсутствуют уравнения некоторых типов, $c, p, s, w_1, w_2, \ldots, w_n \in \Sigma^+$, $Z, W_0, W_1, X_1, X_2, \ldots, X_n, Y_1, Y_2, \ldots, Y_n$ --- переменные. Легко видеть, что, если присутствует уравнение вида $Z = c$, 
то система имеет решение тогда, когда константы $p, s, w_1, w_2, \ldots, w_n$ являются префиксом, 
суффиксом и подсловами константы $c$ соответственно.

Вид уравнений в системе~\ref{eq:system-of-equations} описывает структуру строк, являющихся решениями системы относительно переменной $Z$:
$p$ задает общий префикс, $s$ --- общий суффикс, $w_1, w_2, \ldots, w_n$ --- подслова, которые одновременно встречаются в строке-решении.

Уравнение вида $Z = c$ будем называть константным уравнением, $Z = p W_0$ --- уравнением на префикс, $Z = W_1 s$ --- на суффикс,
$Z = X w Y$ --- на инфикс.


% В этом разделе вводятся базовые понятия, в терминах которых ниже формализуется
% абстрактный анализ строк: стандартный морфизм, трёхпеременная система линейных
% уравнений в словах и пара $(\sigma, S)$, которую мы будем называть
% \emph{строковым свойством}. Семейство строковых свойств с фиксированным
% морфизмом $\sigma$ образует полурешётку, неформально описанную в конце раздела
% (подробное формальное построение --- в разделе~\ref{sec:sp-domain}).

Приведем несколько определений, которые будут использоваться в дальнейшем.

Отметим, что на протяжении всей работы будет рассматриваться фиксированный алфавит $\Sigma$, оснащенный линейным порядком $\leq$ и 
минимальным элементом $\delta_\varnothing$.

\begin{Definition}
Морфизм строк $\sigma \colon \Sigma^* \to \Sigma^*$ называется \emph{алфавитным} (или \emph{неудлиняющим по длине}, \emph{length-non-increasing}), если образ любого символа алфавита $\Sigma$ --- либо символ, либо пустое слово $\varepsilon$.
\end{Definition}

\begin{Definition}[стандартный морфизм]
Будем рассматривать морфизмы $\sigma$, задаваемые разбиением алфавита $\Sigma$ на подмножества $\Sigma_1, \ldots, \Sigma_k, \Sigma_\varepsilon$ при $|\Sigma_i| > 1$, 
которые идемпотентны, не увеличивают букву относительно алфавитного порядка и сохраняют это разбиение, то есть
\[
  \forall \delta \in \Sigma_i \quad
  \sigma(\delta) \leq \delta \;\land\; \sigma(\sigma(\delta)) = \sigma(\delta) \;\land\; \sigma(\delta) \in \Sigma_i,
  \qquad
  \forall \delta \in \Sigma_\varepsilon \quad \sigma(\delta) = \varepsilon .
\]
Такие морфизмы будем называть \emph{стандартными}.
\end{Definition}

В стандартном морфизме для $\delta \in \Sigma_i$ значение $\sigma(\delta)$ --- минимальный относительно $\leq$ символ в $\Sigma_i$.

\begin{Definition}
\label{def:morphorder}
Для стандартных морфизмов $\sigma_1$ и $\sigma_2$ полагаем $\sigma_1 \preceq \sigma_2$ тогда и только тогда, когда $\sigma_2 \circ \sigma_1 = \sigma_2$.
\end{Definition}

Для любого неудлиняющего морфизма $\theta$ с условием $\forall \delta\; \theta(\delta) \preceq \delta$ неподвижная точка $\theta^*$ является стандартным морфизмом относительно некоторого разбиения алфавита.

% \begin{Proposition}
%   Стандартные морфизмы образуют полурешетку с
%   частичным порядком, заданным определением~\ref{def:morphorder}.
  
%   \end{Proposition}
  
%   \begin{proof}
%       Проверим явно аксиомы полурешетки.
%   \begin{enumerate}
%       \item Рефлексивность следует из определения.
%       \item Транзитивность: если $\sigma_1 \circ \sigma_2 = \sigma_1$ и $\sigma_2 \circ \sigma_3 = \sigma_2$, то $\sigma_1 \circ \sigma_3 = \sigma_1$, т.к.
%       $\forall \delta \in \Sigma \; \sigma_1(\sigma_3(\delta)) = \sigma_1(\sigma_2(\sigma_3(\delta))) = \sigma_1(\sigma_2(\delta)) = \sigma_1(\delta)$.
%       \item Антисимметричность: предположим, что $\exists \delta \in \Sigma : \sigma_1(\delta) \neq \sigma_2(\delta)$. 
%       Тогда $\sigma_1(\sigma_1(\delta)) = \sigma_1(\delta) \neq \sigma_1(\delta) = \sigma_1(\sigma_2(\delta))$, противоречие.
%   \end{enumerate}
  
%   % Согласованно с порядком введем операцию $\sqcup = sup$ (будем называть её синонимично обобщением или join для двух стандартных морфизмов).
%   % Поскольку класс стандартных морфизмов рассматривается для одного алфавита, то морфизм однозначно задается разбиением и операция $\sqcup$ задается следующим образом:
%   \end{proof}

\begin{Definition}[Строковое свойство]
\label{def:string-property}
Зафиксируем некоторый алфавит $\Sigma$, будем иметь целью описание некоторых подмножеств моноида $\Sigma^+$.

\textit{Строковым свойством (String Property, SP)} назовем пару $(\sigma, E)$, где 
$\sigma:~\Sigma^*~\to~\Sigma^*$ --- стандартный морфизм, $E$ --- система линейных уравнений в словах над алфавитом $\Sigma' = im(\sigma|_{\Sigma}) \setminus \{\varepsilon\}$
вида, представленного в системе~\ref{eq:system-of-equations}.

Каждое строковое свойство задает в $\Sigma^*$ некоторый язык  $L_{\sigma, E} = \Lang((\sigma, E)) = \sigma^{-1}(\Lang(E))$.  
Здесь $\Lang(E)$ --- $Z$-проекция множества решений системы $E$ как языка над
$\Sigma'$.
\end{Definition}

\begin{Definition}[абстракция и конкретизация]
\label{def:abstraction-concretisation}
Пусть $S$ --- множество конкретных значений данных, а $\langle L^\Delta, \sqcup, \sqcap \rangle$ --- решётка абстрактных значений.
\emph{Абстракцией} называется функция $\alpha \colon 2^S \to L^\Delta$, сопоставляющая множеству конкретных значений обобщающее их абстрактное значение.
\emph{Конкретизацией} называется функция $\gamma \colon L^\Delta \to 2^S$, сопоставляющая абстрактному значению множество описываемых им конкретных значений.

Пара $(\alpha, \gamma)$ задаёт \emph{связь Галуа} (Galois connection), если для любого $A \in 2^S$ выполнено $A \subseteq \gamma(\alpha(A))$ и для любого $\rho \in L^\Delta$ выполнено $\rho \preceq \alpha(\gamma(\rho))$.
\end{Definition}

В случае строкового свойства $(\sigma, E)$ роль множества конкретных значений играет $2^{\Sigma^*}$, а абстрактным значением выступает само строковое свойство; его конкретизация --- задаваемый им язык $L_{\sigma, E}$.

Конкретизацию строкового свойства будем также обозначать
$\gamma\bigl((\sigma, E)\bigr) =  L_{\sigma, E}$.

Зафиксируем стандартный морфизм $\sigma$ и рассмотрим все строковые свойства
$(\sigma, E_1)$ и $(\sigma, E_2)$ с этим морфизмом. Между ними естественно вводится порядок по
включению конкретизаций:
\[
    (\sigma, E_1) \;\preceq\; (\sigma, E_2)
    \;\;\iff\;\;
    \gamma\bigl((\sigma, E_2)\bigr) \,\subseteq\, \gamma\bigl((\sigma, E_1)\bigr).
\]

Совокупность таких свойств с фиксированным $\sigma$ обозначим $\SP(\sigma)$.

Для статического анализа строковых переменнных будем использовать абстрактный домен строковых объектов --- 
объектов, содержащих в себе строковые свойства как компоненты.

Среди стандартных морфизмов выделим два канонических:
\begin{itemize}
    \item тождественный морфизм $\mathrm{id}$: $SP = (\mathrm{id}, E)$ абстрагирует 
    \emph{значение} строки;
    \item морфизм $\sigma_{\forall\to\delta_\varnothing}$, склеивающий весь алфавит
    в выделенный символ $\delta_\varnothing$: $SP = (\sigma_{\forall\to\delta_\varnothing}, E)$,
    задаёт унарное представление нижней границы \emph{длины} строки.
\end{itemize}

\begin{Definition}[абстрактный строковый объект]
\label{def:sob-element}
Зафиксируем алфавит $\Sigma$ и конечный набор стандартных морфизмов
$\sigma_1, \ldots, \sigma_n$.
\emph{Абстрактным строковым объектом} называется кортеж
\[
    \nu \;=\; \bigl\langle\, \mathrm{val}(\nu),\; \mathrm{len}(\nu),\;
        \mathrm{prop}_1(\nu),\, \ldots,\, \mathrm{prop}_n(\nu) \,\bigr\rangle,
\]
в котором $\mathrm{val}(\nu) \in \SP(\mathrm{id})$,
$\mathrm{len}(\nu) \in \SP(\sigma_{\forall\to\delta_\varnothing})$ и
$\mathrm{prop}_i(\nu) \in \SP(\sigma_i)$ для каждого $i = 1, \ldots, n$.
Его конкретизацией называется пересечение конкретизаций всех компонент:
\[
    \gamma(\nu) \;=\; \gamma\bigl(\mathrm{val}(\nu)\bigr) \,\cap\, \gamma\bigl(\mathrm{len}(\nu)\bigr)
        \,\cap\, \bigcap_{i=1}^{n} \gamma\bigl(\mathrm{prop}_i(\nu)\bigr).
\]
\end{Definition}

\begin{Definition}
\label{def:sob}
\emph{Абстрактным доменом строковых объектов} называется редуцированное произведение
полурешёток строковых свойств
\[
    \SOB \;=\; \SP(\mathrm{id}) \treduced \SP(\sigma_{\forall\to\delta_\varnothing})
        \treduced \bigreduced{i=1}{n} \SP(\sigma_i),
\]
$\SP(\mathrm{id})$ --- домен \emph{значения} строки, $\SP(\sigma_{\forall\to\delta_\varnothing})$ ---
домен \emph{длины} (унарное представление), а $\bigreduced{i=1}{n} \SP(\sigma_i)$ ---
необязательное редуцированное произведение остальных свойств, задаваемых стандартными
морфизмами $\sigma_1, \ldots, \sigma_n$.

Конкретизация задаётся пересечением конкретизаций компонент (см.\ определение~\ref{def:sob-element}).
\end{Definition}


\section{Нормализация свойства}

\begin{Definition}
  Словарем будем называть множество строк $D = \{w_1, w_2, \ldots, w_n\}$, обладающее 
  свойством: никакая строка этого множества не является подстрокой другой строки из $D$.
  \end{Definition}
  
  \begin{Definition}
    \label{def:extdict}
  Расширенным словарем будем называть тройку множеств $D = \langle P, S, I \rangle$, 
  где $P$ --- множество префиксов, $S$ --- множество суффиксов, $I$ --- множество инфиксов и при этом:
  \begin{itemize}
      \item $I$ удовлетворяет определению словаря;
      \item строки из $I$ не являются подстроками (возможно, несобственными) строк из $P$ или $S$;
      \item $P$ и $S$ --- пустые множества или множества из одного элемента.
  \end{itemize}
  
  Строки, обозначающие префикс и суффикс расширенного словаря будем сопровождать 
  маркерами начала и конца строки, таким образом они будут иметь вид $\textasciicircum w$ и $w \$$ соответственно.
  Расширенный словарь будем коротко записывать как множество $D~=~\{\textasciicircum p, s \$, w_1, w_2, \ldots, w_n\}$.
  \end{Definition}
  
  Таким образом, расширенный словарь, не имеющий префиксов и суффиксов, является словарем, и вышеобозначенная запись в виде 
  множества строк объединяет определения словаря и расширенного словаря. Поэтому далее под словарем будет
  подразумеваться расширенный словарь, если явно не указано отсутствие ограничения на наличие префиксов и суффиксов.
  
  Условия, отличающие словарь от обычного множества (с выделенным, возможно, суффиксом и префиксом), необходимы,
  чтобы однозначно сопоставить язык, задаваемый словарем, и сам словарь.
  
  \begin{Definition}
  Языком словаря $D$ будем называть множество строк 
  
  $\Lang(D)~=~\{a \in \Sigma^+ \;|\; \forall w_i \in I  \quad a \in p\Sigma^* \, \cap \, \Sigma^*s \, \cap \, \Sigma^*w_i\Sigma^* \}$.
  \end{Definition}

  Многообразие языков, задаваемых рассширенными словарями, представляет собой класс языков, задаваемых системами линейных уравнений вида~\ref{eq:system-of-equations} без 
  уравнений $Z = c$ (при фиксированном алфавите).

  Отметим, что рассматриваемые в данной работе языки входят в некоторый подкласс регулярных языков, что можно использовать
  для обоснования их некоторых свойств.

  \begin{Definition}
    Язык $L \subseteq \Sigma^+$ называется \emph{строго $k$-тестируемым (strictly $k$-testable)},
    если существуют число $k>0$ и четыре множества $U,V,W,F$ такие, что
    \[
    F, U, V \subseteq \Sigma \cup \Sigma^2 \cup \cdots \cup \Sigma^{k-1},
    \qquad
    W \subseteq \Sigma^k,
    \]
    и
    \[
    L=\big[(U\Sigma^* \cap \Sigma^*V)\setminus \Sigma^*W\Sigma^*\big]\cup F.
    \]
    
    Язык, строго $k$-тестируемый для некоторого $k>0$, называется \emph{строго локально тестируемым (strictly locally testable, SLT)}.
    \end{Definition}
    
    \begin{Definition}
    Язык $L \subseteq \Sigma^+$ называется \emph{$k$-тестируемым (k-testable)}, если он является булевой
    комбинацией строго $k$-тестируемых языков.
    
    Язык называется \emph{локально тестируемым (locally testable, LT)}, если он $k$-тестируем для некоторого $k>0$.
    \end{Definition}
    
  
    Интерес представляют языки, задаваемые системой уравнений~\ref{eq:system-of-equations}, а именно языки вида 

    \begin{equation}
    L = \bigcap_{i=1}^n L_i, \text{где } L_i \text{ имеет вид либо } \Sigma^* w_i \Sigma^*, 
    \text{ либо } w_i \Sigma^* \text{ либо } \Sigma^* w_i , w_i \in \Sigma^*.
    \end{equation}
    
    Поскольку язык вида $\Sigma^* \setminus \Sigma^* w_i \Sigma^*$ локально тестируем (более того, строго локально тестируем),
    то $\Sigma^* w_i \Sigma^*$ также локально тестируем. Очевидно, что $w_i \Sigma^*$ и $\Sigma^* w_i$ также строго локально тестируемы,
    из чего следует, что $L$, представленный как пересечение этих компонент --- LT-язык.
    
    Проекция решения на переменную $Z$ позволяет кодировать некоторое множество строк, поскольку из системы известны
    ограничения на их префиксы, суффиксы и инфиксы. Несмотря на то, что $Z$-проекция решения --- регулярный язык,
    при возрастании числа уравнений возрастает трудоемкость построения для него альтернативного представления, поэтому
    хранение словаря --- набора строк из уравнений системы --- является наиболее эффективным. 
    
    Например, пустое множество решений может задаваться любой системой с набором противоречащих друг другу уравнений,
    что легко проверяется анализом константных строк, входящих в уравнения, но трудоемко, если использовать ДКА для каждого уравнения и их пересечения.

    Заметим, что представление языка в виде словаря с присоединеной альтернативой в виде единственной строки эквивалентна представлению
    в виде системы уравнений вида~\ref{eq:system-of-equations} при фиксированном алфавите.
    
\begin{Definition}[Задача внутренней нормализации строкового свойства]
Представление языка $L \subseteq \Sigma^*$ $L = \bigcap_{i=1}^n L_i$ ($L_i$ имеет вид $ \Sigma^* w_i \Sigma^*$ или $w_i \Sigma^*$ или $\Sigma^* w_i$ или $w_i$, $w_i \in \Sigma^*$)
является нормализованным, если выполняются следующие условия.
\begin{enumerate}
\item Cреди компонент пересечения нет такой $L_i$, что $L_i \subseteq L_j$ для $i \neq j$ (отсутствие избыточности). 
\item Не существует такого $L_{n+1}$, что $L = L \cap L_{n+1}$ и при этом не нарушено предыдущее условие (полнота представления).
\end{enumerate}

Соответственно, для нормализации представления необходимо убрать компоненты, нарушающие первое условие, и добавить компоненты,
нарушающие второе.
\end{Definition}

Отметим, что проверка отсутствия избыточности представления более тривиальна, чем поиск полноты представления.

Если язык является множеством-синглетоном, то нормализация избыточности его представления сводится к проверке остальных уравнений на 
подстроки этого элемента, иначе система противоречива, и язык эквивалентен пустому множеству. Иначе говоря, условия накладываемые 
на расширенный словарь в определении~\ref{def:extdict} обеспечивают отсутствие избыточности представления.
% TODO: можно развить обоснование, но пока и так сойдет

Для внутренней нормализации строкового свойства необходимо следующее утверждение, представленное в статье~\cite{NepeivodoaAfanasyev2025} в терминах неизбегаемых подстрок.
Приведем здесь переформулировку с альтернативным обоснованием, проиллюстрировав связь с известными свойствами класса LT-языков.

Будем использовать теорему 7~\cite{Zalcstein1972} о строении языка $L_u = \Sigma^* u \Sigma^*$ и ходом ее доказательства.

\begin{Theorem}[Зальстейн, теорема 7]
\label{thm:zalcstein-7}
Пусть $L_u = \Sigma^* u \Sigma^*$, $u \in \Sigma^+$. Следующие условия эквивалентны:
\begin{enumerate}
    \item Либо $|\Sigma| \geq 3$, либо $\Sigma = \{a, b\}$ и $u$ не имеет вида $a^n b$ или $a b^n$ ни для какого $n \geq 1$.
    \item Минимальный ДКА $\mathcal{A}(L_u)$, распознающий $L_u$, является \emph{почти сильно связным} ($0$-транзитивным): у него есть поглощающее (и единственное допускающее) состояние $d$, и для любых двух состояний $q', q'' \in Q \setminus \{d\}$ найдётся такое слово $v \in \Sigma^*$, что $q' \cdot v = q''$.
    \item Синтаксическая полугруппа $\mathcal{S}(L_u)$ --- нильпотентное расширение конечного $0$-простого комбинаторного идеала.
\end{enumerate}
\end{Theorem}

Будем использовать обозначение $u[i:j]$ для подстроки $u$ от индекса $i$ до индекса $j$ включительно (первый символ строки имеет индекс 1).

В ходе доказательства теоремы~\ref{thm:zalcstein-7} используется следующая лемма.

\begin{Lemma}[Зальстейн, лемма 2]
\label{lem:zalcstein-2}
Пусть $u = u_1 u_2 \ldots u_m \in \Sigma^+$, $L_u = \Sigma^* u \Sigma^*$ как в теореме~\ref{thm:zalcstein-7}, и $\mathcal{A}(L_u)$ --- минимальный ДКА, 
распознающий $L_u$, с множеством состояний $Q = \{q_0, q_1, \ldots, q_m\}$, 
причём $q_i = q_0 u[1:i]$ для $1 \leq i \leq m$, а $q_m = d$ --- 
поглощающее допускающее состояние. Тогда для всякого состояния $q_i$ ($i < m$) и
 любого слова $v \in \Sigma^+$ верно:
\begin{itemize}
    \item $q_i v = d$ тогда и только тогда, когда $u$ является подсловом строки $u[1:i] v$;
    \item в противном случае $q_i  v = q_j$, где $u[1:j]$ --- наибольший по длине префикс $u$, являющийся суффиксом строки $u[1:i] v$.
\end{itemize}
\end{Lemma}

Будем использовать вышеприведенные утверждения для обоснования следующего утверждения о полноте представления.


\begin{Proposition}
\label{prop:fullness-of-representation}
Пусть $L$ --- язык вида, представленного в задаче нормализации строкового свойства, в котором выполнено условие отсутствия избыточности.
Тогда,
\begin{itemize}
\item если $|\Sigma| > 2$, свойство полноты представления всегда выполняется;
\item если $|\Sigma| = 2$, свойство полноты представления могут нарушать только компоненты вида $\Sigma^* ab^n \Sigma^*$ и $\Sigma^* b^n a \Sigma^*$
($\Sigma = \{a, b\}$).
\end{itemize}
\end{Proposition}

\begin{proof}

0. Отметим, что компоненты вида $w_i \Sigma^*$, $\Sigma^* w_i$ и $w_i$ не могут нарушать свойство полноты представления.

1. Пусть строковое свойство задает язык $L = \bigcap_{i=1}^n \Sigma^* w_i \Sigma^* \;  \bigcap \; p \Sigma^* \; \bigcap \; \Sigma^* s$
(здесь $p$ и/или $s$, возможно, пустые строки, т.е. ограничений на префикс и суффикс может не быть). 
Предположим, что $u$ --- некоторый фактор, не являющийся подстрокой ни в $p$, ни в $s$, ни в $w_i$ для любого $i$, что
$L_u = \Sigma^* u \Sigma^*$ нарушает свойство полноты представления.

1.1. Рассмотрим случай, когда $L_u$ --- как в теореме~\ref{thm:zalcstein-7}. 
По предположению, любое слово $w$ языка $L$ должно быть принято $\mathcal{A}(L_u)$.

Согласно лемме~\ref{lem:zalcstein-2}, которая используется в статье для доказательства теоремы~\ref{thm:zalcstein-7}, для 
$ v \in \Sigma^+  \; q_i v = d \iff $
 $u$ --- подстрока $u[1:i]v$. Поскольку изначально $u$ не является подстрокой ни в $p$, ни в $s$, ни в $w_i$,
 из леммыы получаем, что $s$, $p$ и все $w_i$ по отдельности не могут быть приняты ДКА $\mathcal{A}(L_u)$.

  Из теоремы~\ref{thm:zalcstein-7} 
 следует, что $\exists v_j : q_j v_j = q_0, \; j \leq m-1$. Так как $q_0 p \neq d$,\; $q_0 s \neq d$ и $q_0 w_i \neq d$, для всех $i$,
 $\exists v_p, v_1, \ldots, v_n :\; q_0 p v_p = q_0,\; q_0 w_i v_i = q_0, \; i \leq n$.

 Это отражение того, что автомат связный, если убрать допускающее (поглощающее) состояние $d$. Т.е. 
 пока мы не прочитаем $u$ как подстроку полностью, мы можем некоторым образом вернуться в начальное состояние.

 Сконструируем слово 

 $$
 w = p v_p w_1 v_1 \ldots w_n v_n s.
 $$

 По конструкции видно, что $w \in L$ и по предположению оно принимается $\mathcal{A}(L_u)$.
 
 Но $q_0 w = q_0 \,p v_p \, w_1 v_1 \ldots w_n v_n\, s = q_0 \,w_1 v_1 \ldots w_n v_n\, s = \ldots  = q_0 \,w_n v_n \,s = q_0 \, s$,
и при этом $q_0\, s \neq d$. Получаем противоречие, из чего следует вывод, что языки, описанные в теореме~\ref{thm:zalcstein-7}, не могут 
нарушать свойство полноты представления.

1.2. Пусть $\Sigma = \{a, b\}$.

1.2.A. Пусть $u = a b^m$. Тогда если найдутся $w_i, w_j$, удовлетворяющие регулярному выражению
$b^*a(a|b)^+b^m$ и не содержащие $u$ как подстроку, то $u$ обязательно содержится во всех словах $L$.

1.2.Б. Пусть $u = a^m b$. Тогда если найдутся $w_i, w_j$, удовлетворяющие регулярному выражению
$a^*b(a|b)^*a^m$ и не содержащие $u$ как подстроку, то $u$ обязательно содержится во всех словах $L$.
 
С точностью до переименования символов, получаем описания некоторых случаев, когда $u$ определенного утверждением
вида действительно нарушает свойство полноты представления и может быть добавлено в представление языка $L$.

\end{proof}

Случаи, когда компонента вида $\Sigma^*a b^m \Sigma^*$ (или $\Sigma^*a^m b \Sigma^*$) в двухбуквенном алфавите
может быть добавлена в представление языка $L$, описаны в доказательстве утверждения~\ref{prop:fullness-of-representation} довольно
широко, но не полностью, поскольку не рассмотрены варианты, когда префикс или суффикс (возможно, совместно с некоторой подстрокой)
вызывают необходимость пополнения представления языка $L$.

\begin{Example}
Пусть $L = \{aa\Sigma^*b\}$ и $\Sigma = \{a, b\}$. Тогда $aab$ также содержится во всех словах $L$.
\end{Example}

\begin{Example}
Пусть $L = \Sigma^* abbb \Sigma^*\, \bigcap\, \Sigma^* baa \Sigma^*\, \bigcap\, \Sigma^* ab $, $\Sigma = \{a, b\}$.
Тогда слова $bbba$ и  $aab$ также содержатся во всех словах $L$.
\end{Example}


\section{LTL-свойства}

Необходимо различать язык, который задает система уравнений сама по себе и язык, который задается строковым свойством.
Как было уже отмечено ранее, язык строкового свойства задается как $\Lang((\sigma, E)) = \sigma^{-1}(\Lang(E))$.
В системе уравнений используется ограниченный алфавит, поэтому он описывает лишь проекцию языка свойства на множество 
строк в алфавите $\Sigma' = im(\sigma|_{\Sigma}) \setminus \{\varepsilon\}$.

Для тождественного отображения язык строкового свойства совпадает с языком его системы уравнений --- \textit{значения} строкового объектв.
Как будет показано далее, локально тестируемые языки удобны для анализа, поэтому имеет место вопрос, задает ли строковое свойство LT язык.
В этом случае строковое свойство действует в некотором смысле как система уравнений (поскольку сами системы задают LT языки).

Для дальнейшего рассуждения будем использовать теорему о характеризации LT-языков~\cite{McNaughton1974,BrzozowskiSimon1973,Caron2000}.

\begin{Theorem}[Caron, теорема 2.5]
\label{thm:caron-lt}
Пусть $L \subseteq \Sigma^+$ --- регулярный язык. Следующие условия эквивалентны:
\begin{enumerate}
    \item $L$ --- локально тестируемый (LT) язык.
    \item Синтаксическая полугруппа $\mathcal{S}(L)$ языка $L$ локально идемпотентна и коммутативна, то
     есть для всякого идемпотента $e \in E(\mathcal{S}(L))$ подполугруппа $e\mathcal{S}(L)e$ идемпотентна и коммутативна.
\end{enumerate}
\end{Theorem}


\begin{Proposition}
Строковое свойство $(\sigma, E)$ задает LT язык тогда и только тогда, когда верно одно из следующего:
\begin{itemize}
% \item $S$ не имеет уравнений;
\item $|im(\sigma|_{\Sigma})| = 1$;
\item $im(\sigma|_{\Sigma}) = \{\alpha, \varepsilon\}$ и $E$ содержит только уравнения вида $Z = X \alpha Y$, $Z = X \alpha$, $Z = \alpha Y$;
\item $\sigma$ --- нестирающий морфизм;
\item $\sigma$ --- стирающий морфизм, $E$ --- система уравнений, имеющая только уравнения вида 
$Z = X_i \alpha_i Y_i$, где $\alpha_i \in im(\sigma|_{\Sigma})$.
\end{itemize}
\end{Proposition}

\begin{proof}
    Отметим, что для случая $|im(\sigma)| = 1$ свойство задает LT язык, поскольку в этом случае $im(\sigma|_{\Sigma}) = \{\alpha\}$ или $im(\sigma|_{\Sigma}) = \{\varepsilon\}$.
    Вторая альтернатива тривиальна, а в первом случае это свойство задает ограничение лишь на длину строки.

    \underline{1.} Пусть $\sigma$ --- нестирающий морфизм. 

    \underline{1.1.} Если $E = \{Z = c\}$, то, поскольку алфавит конечный, $\Lang((\sigma, E)) = \sigma^{-1}(c) = \bigcup_{i=1}^n \{c_i\}, c_i \in \Sigma^+$, что дает LT язык.

    \underline{1.2.} Система уравнений не представляет собой константное уравнение. 
    Зададим отображение уранения во множество строк, которое оно задает в данном строковом свойстве.
    \begin{align*}
        Z = X w Y &\mapsto \Sigma^* \sigma^{-1}(w) \Sigma^* \\
        Z = X w &\mapsto \Sigma^* \sigma^{-1}(w) \\
        Z = w Y &\mapsto \sigma^{-1}(w) \Sigma^*
    \end{align*}

    % (Прим.: напомним, что $\sigma^{-1}(w) = \{v \in \Sigma^* \mid \sigma(v) = w\}$ и запись $A B C$, где $A, B, C$ --- множества строк,
    % означает множество $\{abc | a \in A, b \in B, c \in C\}$)

    Очевидно, что множество строк, задаваемых уравнением каждого вида --- LT язык, а следовательно их пересечение также является LT языком.

    \underline{2.} Пусть $\sigma$ --- стирающий морфизм, $im(\sigma|_{\Sigma}) \neq \{\varepsilon\}$.

    Для доказательства далее будем использовать алгебраическую характеризацию LT языков (теорема~\ref{thm:caron-lt}):

    Будем обозначать $[w]$ --- элемент синтаксической полугруппы $\mathcal{S}(L)$, соответствующий строке $w \in \Sigma^+$. 
    
    Поскольку язык задается строковым свойством, то определение синтаксической конгруэнции $\sim_L$ можно эквивалентно переписать так:
    \begin{center}
    $u \sim_L v$ если $ \forall x, y \in \Sigma^+ \; \sigma(xuv) \text{ удовлетвояет системе } E \iff \sigma(xvy) \text{ удовлетвояет системе } E$.
    \end{center}

    Пусть $\sigma(\gamma) = \varepsilon$. Тогда $[\gamma]$ --- идемпотент в $\mathcal{S}(L)$.

    \underline{2.1.} Покажем, что случай $E = \{Z = c\}$ не дает LT язык. $[\gamma c \gamma]$ не является идемпотентом в $[\gamma] \mathcal{S}(L) [\gamma]$, поскольку
    $\gamma c \gamma \not \sim_L \gamma c \gamma \gamma c \gamma$.

    \underline{2.2.} Иначе расматриваем неконстантные системы уравнений. 

    Пусть $A$ --- множество символов, встречающихся в строках системы уравнений, 
    $A \subseteq im(\sigma|_{\Sigma}) \setminus \{\varepsilon\} \subseteq \Sigma$.

    \underline{2.2.1.} $|A| = 1, A = \{\alpha\}$.
    По предположению о локальной тестируемости языка, верно следующее:
    \begin{align*}
        [\gamma \alpha \gamma] [\gamma \alpha \gamma] &= [\gamma \alpha \gamma] \\
        [\gamma \alpha \alpha \gamma] [\gamma \alpha \alpha \gamma] &= [\gamma \alpha \alpha \gamma]
    \end{align*}

    Тогда для натурального $n$ $\alpha^n \in L$ и $\alpha \sim_L \alpha^n$, отсюда следует, что ограничения могут быть только однобуквенными.
    Тогда система уравнений имеет вид $\{Z = \alpha X_1; Z = X_2 \alpha\}$, или $\{Z = X \alpha\}$, или  $\{Z = \alpha X\}$, или 
    $\{Z = X \alpha Y\}$.

    \begin{enumerate}
        \item $\{Z = X \alpha Y\}$: является LT языком, а именно $\Sigma^* \sigma^{-1}(\alpha) \Sigma^*$.
        \item $\{Z = \alpha X_1; Z = X_2 \alpha\}$: для $im(\sigma|_{\Sigma}) = \{\alpha, \varepsilon\}, \gamma \mapsto \varepsilon$ язык является LT (а именно 
        $(\sigma^{-1}(\varepsilon) \sigma^{-1}(\alpha) \sigma^{-1}(\varepsilon))^+ = A^+ \setminus \sigma^{-1}(\varepsilon)$); если $\exists \beta \mapsto \beta$, 
        язык не является LT (контрпример: слова $\gamma^k \alpha \gamma^k  \beta \gamma^k$ и
        $\gamma^k \alpha \gamma^k \beta \gamma^k \alpha \gamma^k$ имеют одинаковые $k$-подстроки, но первое не приналежит языку).
        \item $\{Z = \alpha X\}$ рассуждение аналогично предыдущему пункту.
        \item $\{Z = X \alpha\}$: полностью симметрично предыдущему пункту. 
    \end{enumerate}

    \underline{2.2.2.} $|A| > 1, A = \{\alpha_1, \ldots, \alpha_n\}, \sigma(\alpha_i) = \alpha_i$.

    Аналогично рассуждению выше, из комммутативности и идемпотентности 
    $$\alpha_i\alpha_i \sim_L \alpha_i, \quad \alpha_i\alpha_j \sim_L \alpha_j\alpha_i$$.
    
    Пусть $w = \alpha_{k_1} \ldots \alpha_{k_m} \in A^+,\; k_j \in \{1, \ldots, n\}$, $w$ --- минимальная по длине строка, удовлетворяющая системе уравнений (minimal common superstring).
    Тогда $w~\sim_L~\alpha_{i_1}~\alpha_{i_2}~\ldots~\alpha_{i_m}$, $(i_1, \ldots, i_m)$ --- такая 
    перестановка множества $\{k_1, \ldots, k_m\}$, где если $\alpha_{i_j} = \alpha_{i_{(j + s)}}$, то $\forall t \leq s, t \geq 0 \; \alpha_{i_{(j + t)}} = \alpha_{i_j}$.

    Воспользовавшись свойством идемпотентности для соседних символов, получаем $w \sim_L w' = \alpha_{p_1} \alpha_{p_2} \ldots \alpha_{p_r}$, $\alpha_{p_i}$ --- различные символы, $r \leq m$.

    \underline{2.2.2.А.} Покажем, что для случая подстрок длины $\geq 2$ в оставшемя подслучае язык не является LT.

    Если в системе уравнений есть уравнения с подстроками (с префиксами или суффиксами) длины 2 или более, то верно одно из двух:
    \begin{enumerate}
        \item $|w'| < |w|$ и, так как $w$ --- минимальная строка, $w'$ не удовлетворяет системе уранений, т.е. $w' \not \sim_L w$.
        \item $|w'| = |w|$, $w' \sim_L w$, $w'$ --- одна из минимальных строк, удовлетворяющих системе уравнений. \\
        Слово $w'$ содержит только различные символы, значит в уравнении системы с подстрокой длины $\geq 2$ будут использоваться различные символы, условно $\alpha_a, \alpha_b$.
    Сконструируем слово $w''$, полученное из $w'$ перестановкой $\alpha_a$ и $\alpha_b$. Поскольку сочетание $\alpha_a \alpha_b$ теперь не встречается в $w'$, то 
    получаем противоречие: $w' \not \sim_L w$. 
    \end{enumerate}

    

    \underline{2.2.2.Б.} Покажем утверждение теоремы для систем, содержащих уравнения только вида $Z = X \alpha_i Y$, $Z = X \alpha_j$, $Z = \alpha_k Y$.
    
    Заметим, что наличие уравнений на префикс или суффикс выводит язык из класса LT. Пусть есть уравнение $Z = \alpha_{p_1} W$, и $w'$ удовлетворяет системе уравнений.
    По построению $\alpha_{p_2} \neq \alpha_{p_1}$, значит $w'' = \alpha_{p_2} \ldots \alpha_{p_r} \alpha_{p_1} \not \sim_L w'$. Симметричное расссуждение проводится и для суффикса.

    Для систем, состоящих только из уравнений вида $Z = X \alpha Y$, легко видеть, что задаваемый их строковым свойством язык является LT.
    \begin{align*}
    \Lang((\sigma, \{Z = X \alpha Y\})) = \Sigma^* \sigma^{-1}(\alpha) \Sigma^* \\
    L = \Lang((\sigma, E)) =  \bigcap_{i=1}^k \Sigma^* \sigma^{-1}(\alpha_{p_i}) \Sigma^* 
    \end{align*}
    
    Таким образом, покрыв рассуждением все языки, задаваемые строковыми свойствами, получаем утверждение теоремы.

\end{proof}

% \begin{Proposition}


\section{Стандартные и монотонные морфизмы}

Зададим следующие операции на классе стандартных морфизмов:

\begin{align}
    \label{morphmeet}
    (f \sqcap g)(x) &= \min(f^{-1}(f(x)) \cap g^{-1}(g(x))) \quad \text{ (meet)}; \\
    \label{morphjoin}
    (f \sqcup g)(x) &= \min [x]_{f,g} \quad \text{ (join)}.
\end{align}

Здесь $[x]_{f,g}$ --- класс эквивалентности отношения $\sim_{f, g}$, которое является транзитивным замыканием отношения на символах алфавита 
для фиксированной пары морфизмов $f, g$ следующим образом:

\begin{equation}
x \sim f(x) \wedge x \sim g(x) \text{ для всех } x.
\end{equation}

\begin{Proposition}
Операции join и meet (соотношения~\ref{morphmeet} и~\ref{morphjoin}) задают решетку на стандартных морфизмах.
Эти операции согласованы с частичным порядком из определения~\ref{def:morphorder}.
\end{Proposition}

\begin{proof}

    Покажем корректность операций относительно порядка.

    1.1. $f \sqcap g \preceq f$. Действительно, $(f \sqcap g)(x) \in f^{-1}(f(x))$, значит $f((f \sqcap g)(x)) = f(x)$ для всех $x$.

    1.2. $g \preceq f \iff f \sqcap g = g$. 

    ($\Rightarrow$). Пусть $f(g(x)) = f(x)$, тогда, т.к. $g(g^{-1}(g(x))) = g(x)$, то $f(g(g^{-1}(g(x)))) = f(g^{-1}(g(x))) = f(g(x)) = f(x)$, т.е.
    $f^{-1}(f(x)) \supseteq g^{-1}(g(x))$ (это соответствует тому, что разбиение алфавита в меньшем морфизме всегда вкладывается в разбиение алфавита в большем морфизме).

    Отсюда следует, что $f^{-1}(f(x)) \cap g^{-1}(g(x)) = g^{-1}(g(x))$, что доказывает первую часть утверждения.

    ($\Leftarrow$). Из определения операции получаем вложение $f^{-1}(f(x)) \supseteq g^{-1}(g(x))$ (иначе для элементов вне пересечения поведение meet-морфизма отличалось бы от $g$).
    Далее, $f(g(g^{-1}(g(x)))) = f(x)$, т.е. $f(g(g(x))) = f(g(x)) = f(x)$.

    2.1. $f \sqcup g \succeq f$. Действительно, $f(x) \in [x]$, значит $(f \sqcup g)(f(x)) = (f \sqcup g)(x)$ для всех $x$.

    2.2. $g \preceq f \iff f \sqcup g = f$. 

    ($\Rightarrow$). Пусть $f(g(x)) = f(x)$. Выше обосновано, что $f^{-1}(f(x)) \supseteq g^{-1}(g(x))$, поэтому $[x] = f^{-1}(f(x))$, значит $\min[x] = f(x)$.

    ($\Leftarrow$). Равенство означает, что $[x]_{f,g} = f^{-1}(f(x))$, значит $f(x) \sim_{f,g} x \sim_{f_g} y \in g^{-1}(g(x))$, что снова приводит в вложению разбиений, и
    аналогично доказательству для meet снова приводит к требуемому утверждению. 

    3. Идемпотентность, коммутативность, ассоциативность обеих операций следует из их определения и определения стандратного морфизма.

    4. Поскольку уже показана согласованность пордка с операциями, законы поглощения обосновать несложно.
    Поскольку $f \succeq (f \sqcap g) $, то $f \sqcup (f \sqcap g) = f$. Полностью по аналогии $f \sqcap (f \sqcup g) = f$.
\end{proof}

% монотонные

Рассмотрим класс монотонных стандартных нестирающих морфизмов. А именно,
зафиксировав алфавит $\Sigma$, снабженный линейным порядком, монотонный морфизм ---
морфизм, удовлетворяющий свойству $\forall a, b \in \Sigma, a \leq b \implies \sigma(a) \leq \sigma(b)$.

Поскольку $\Sigma$ --- линено упорядоченное множество, то верно и следующее: если $\sigma(a) < \sigma(b)$, то $a < b$ для всех $a, b \in \Sigma$.

Монотонные стандартные нестирающие морфизмы не имеют <<разорванных>> разбиений, т.е. для любого $\Sigma_i$
все элементы алфавита в диапазоне (согласно линейному порядку) от $\min(\Sigma_i)$ до $\max(\Sigma_i)$ также входят в $\Sigma_i$. 

\begin{Proposition}
    Порядок на стандартных морфизмах $\preceq$, заданный ранее ($\sigma \preceq \tau \iff \tau \circ \sigma = \tau$), для 
    множества монотонных стандартных нестирающих морфизмов эквивалентен порядку 
    $\preceq_M$, заданному следующим образом:
    $$
    \sigma \preceq_M \tau \iff \forall x \in \Sigma, \tau(x) \preceq \sigma(x).
    $$
\end{Proposition}

\begin{proof}
    Для доказательства приведем явную проверку эквивалентности 
    $$
    \tau \circ \sigma = \tau \iff \forall x \in \Sigma, \; \tau(x) \leq \sigma(x).
    $$

    \underline{$\implies$.} 
    $\tau(\sigma(x)) = \tau(x)$ для всех $x \in \Sigma$. По определению стандартного морфизма $\tau(x) \leq x$, подставляя аргумент ---
    $\tau(x) = \tau(\sigma(x)) \leq \sigma(x)$.
    
    \underline{$\impliedby$.} 
    $\tau(x) = \tau(\tau(x)) \leq \tau(\sigma(x))$ с одной стороны. С другой: $\sigma(x) \leq x$, значит $\tau(\sigma(x)) \leq \tau(x)$. Поскольку
    это верно для всех $x \in \Sigma$, получаем требуемое равенство.
\end{proof}

Добавление монотонности позволяет упростить вычисление операций join и meet для монотонных стандартных нестирающих морфизмов следующим образом:

\begin{align}
  \label{morphmeetmono}
  \bigl(\sigma \sqcap \theta \bigr)(x) = \max (\sigma (x), \theta(x))\\
  \label{morphjoinmono}
  \bigl(\sigma \sqcup \theta\bigr)(x) = (\sigma \circ \theta)^* (x) = (\theta \circ \sigma)^* (x),
\end{align}
здесь $\tau = (\sigma \circ \theta)^*$ --- неподвижная (фиксированнная) точка отображения $\sigma \circ \theta$. 

Эти соотношения --- монотонные аналоги соотношений~\ref{morphmeet} и~\ref{morphjoin}. 

Для конкретного аргумента $x$ элементы разбиения двух морфизмов $\sigma^{-1}(\sigma(x))$ и $\theta^{-1}(\theta(x))$ непрерывны,
пересекаются по непрерывному отрезку и при этом $\min \sigma^{-1}(\sigma(x)) = \sigma(x)$ (что верно и для $\theta$). Поэтому минимальный элемент в пересечении прообразов ---
именно $\max (\sigma(x), \theta(x))$, что и задает операцию meet.

Для монотонных морфизмов класс эквиватентности $[x]_{f,g}$ содержит каждую точку на пути стабилизации отображения 
$x \mapsto  \theta (x) \mapsto (\sigma \circ \theta) (x) \mapsto (\theta \circ \sigma \circ \theta) (x) \mapsto (\sigma \circ \theta \circ \sigma \circ \theta) (x) \mapsto \hdots \mapsto \tau(x)$
и каждый переход этой цепочки невозрастающий. Это прямое следствие монотонности: если $\sigma(x) \leq \theta(x)$, то $\theta(\sigma(x)) \leq \theta(x)$ и 
при этом $\sigma(\theta(x)) \leq \sigma(x)$ (т.к. по определению стандартного морфизма $\theta(x) \leq x$); если же $\theta(x) \leq \sigma(x)$, рассуждаем симметрично. 

Именно поэтому композиции $\sigma \circ \theta$ и $\theta \circ \sigma$ стабилизируются в одной точке для монотонных морфизмов. Для немонотонных морфизмов
это неверно: утверждение $\sigma(x) \leq \theta(x) \implies \theta(\sigma(x)) \leq \theta(x)$ может не выполняться, что приведет к тому, 
что фиксированнные точки композиций будут различаться, что в принципе нарушит условие присутствия точной верхней грани.

\begin{Example}
Пусть $\Sigma = \{a, b, c\}$ с линейным порядком $a < b < c$, и стандартные морфизмы $\sigma, \theta$ заданы следующим образом:
$$
\sigma(a) = a, \sigma(b) = b, \sigma(c) = a, \\
\theta(a) = a, \theta(b) = b, \theta(c) = b.
$$
Тогда $\sigma \circ \theta = \theta$ и $\theta \circ \sigma = \sigma$. 

Таким образом, поскольку эти морфизмы не монотонны, $(\sigma \circ \theta)^*$ и $(\theta \circ \sigma)^*$ будут различаться и, более того,
не будут являться верхними гранями для $\sigma$ и $\theta$, поскольку они несравнимы.

Прямой проверкой также можно убедиться, что определение монотонного meet~\ref{morphmeetmono} для этой пары морфизмов также не даст точную нижнюю грань.

\end{Example}


\section{Кросс-нормализация}
\label{sec:cross-normalization}

Строковые свойства --- компоненты одного строкового объекта, рассматриваемые одновременно. Поэтому в явном виде
множество строк, описываемых строковым объектом, можно описать пересечением 
языков всех его строковых свойств.

\begin{Definition}
Кросс-нормализацией строкового свойства $SP_1$ относительно другого строкового свойства $SP_2$ называется 
добавление всех уравнений допустимого $SP_1$ вида, которые описывают язык $\Lang(SP_1) \cap \Lang(SP_2)$. 
\end{Definition}

Каждое уравнение $SP_2$ описывает язык $I_2$ одного из видов: $\Sigma^* \sigma_2^{-1}(w) \Sigma^*$,
$\Sigma^* \sigma_2^{-1}(w)$, $\sigma_2^{-1}(w) \Sigma^*$, $\sigma_2^{-1}(w)$.

Кросс-нормализация $SP_1$ относительно $SP_2$ решает задачу нахождения такого множества строк $I_1 \subseteq \Sigma^+$, что
$\sigma_1(I_1) \supseteq \sigma_1(I_2)$, при этом $\sigma_1(I_1)$ является множеством строк, задаваемым уравнением (или уравнениями) одного из
допустимых видов.

Неформально, это означает, что в систему уравнений $SP_1$ добавляются все возможные уравнения, которые также 
описывают пересечение языков $SP_1$ и $SP_2$, поскольку строковые свойства для одного объекта имеют целью описать 
множество конкретизации объекта в разных <<проекциях>>.

Из определения строкового свойства $I_2$ уравнения представляет собой объединение тривиальных компонент, тем временем $\sigma_1(I_1)$
должно иметь форму пересечения некоторых также тривиальных компонент. Тривиальные компоненты --- левые, правые и двусторонние идеалы $\Sigma^+$, 
а также множества-синглтоны из одной строки $\{w\}$. Кроме того, т.к. $I_2$ --- язык одного конкретного уравнения относительно морфизма $\sigma_2$,
он является либо левым, либо правым, либо двусторонним идеалом, либо конечным множеством строк. Покрывающее множество $I_1$
должно являться более объемлющим; к примеру, $\sigma_1(I_1)$ не может быть левым идеалом, если $I_2$ --- язык уравнения на инфикс.

\begin{Proposition}
Если $I_2$ --- язык одного уравнения (на префикс, на суффикс, на инфикс, на константу) относительно морфизма $\sigma_2$ и 
$P_1, P_2$ --- некоторые множества строк, такие что $\sigma_1(P_1) \supseteq \sigma_1(I_2)$ и $\sigma_1(P_2) \supseteq \sigma_1(I_2)$, то 
$\sigma_1(P_1 \cap P_2) \supseteq \sigma_1(I_2)$.
\end{Proposition}

Это утверждение не треебует доказательства, поскольку заключается в простом пересечении двух множеств без привязки к свойствам языков.

Поскольку каждое уравнение использует конечную строку (условно, некоторую $u$), покрывающих множеств тривиального вида (синглетоны, левые, правые и двусторонние идеалы)
для $\sigma_1(I_2)$ конечное число, соответствующее не более чем числу всех подстрок в строках множества $\sigma_1(\sigma_2^{-1}(u))$. Отсюда вывод:
минимальный $\sigma_1(I_1)$ является пересечением некоторых тривиальных компонент, каждая из которых задается некоторой подстрокой строки из множества
 $\sigma_1(\sigma_2^{-1}(u))$.


Некоторая подстрока $w$ будет задавать покрывающий идеал или константное уравнение только если выполняется следующее 
\underline{необходимое условие кросс-переноса для LTL-свойств} (в $SP_1$ из $SP_2$):
\begin{equation}
    \sigma_1(\sigma_2^{-1}(w)) \text{ --- синглтон. }
\end{equation}

Это верно, поскольку для LTL-свойств $\sigma_2^{-1}(w)$ --- конечное множество строк, таким образом уравнение может быть добавлено в общее пересечение, 
иначе необходимо взять более короткую подстроку, которая будет гарантированно встречаться во всех строках множества $\sigma_1(\sigma_2^{-1}(w))$.

Поскольку для строковых свойств, задающих LT языки, верно, что неконстантные уравнения задают либо левый, либо правый, либо двусторонний идеал, то и в обратную сторону,
нахождение покрывающего идеала --- достаточно простая задача. Однако это неверно для не-LTL свойств: каждое уравнение такого свойства задает регулярный, но не LT язык, что 
затрудняет нахождение нетривиального языка, представимого уравнением другого строкового свойства. 

Более того, поскольку из LTL-свойств переносятся только некоторые сохраняемые относительно морфизмов подстроки, это 
позволяет передать максимум информации в сравнимые свойства. При этом в несравнимое с данным строковое свойство переносятся
уравнения, которые также возможно перенести в некоторое сравнимое свойство. Для не-LTL свойств подобный механизм не работает
из-за более сложной структуры языков. 

Пусть в некотором уравнении строкового свойства используются только символы из $A \subseteq \Sigma$, 
тогда для не-LTL свойства в худшем случае найдется порядка  $2^{|A|}$ несравнимых по морфизму строковых свойств (поскольку, опираясь на структуру уравнения, различные символы уравнения
могут по-разному быть объединены в разбиении, что существенно влияет на структуру переносимых уравнений). Это 
отражено на рисунке~\ref{fig:nonltl}: на пунктирных стрелках подписаны подстроки, переносимые в несравнимое свойство.

\begin{figure}[htb]
    
\begin{center}
\begin{tikzpicture}[
    y=1.38cm,
    lat/.style={
      rectangle,
      draw,
      thick,
      rounded corners=2pt,
      minimum width=2.3cm,
      minimum height=0.95cm,
      align=center,
      inner sep=4pt,
      fill=white
    },
    lat-empty/.style={
      rectangle,
      draw=black!35,
      thick,
      rounded corners=2pt,
      minimum width=1.65cm,
      minimum height=0.72cm,
      align=center,
      inner sep=2pt,
      fill=gray!14,
      text=black!45
    },
    arrlbl/.style={
      fill=white,
      inner sep=2pt,
      font=\footnotesize,
      pos=0.6,
      above
    },
    tag/.style={
      rectangle,
      draw,
      thick,
      rounded corners=1pt,
      font=\scriptsize,
      inner sep=3pt,
      align=center,
      fill=white
    }
  ]
  \node[lat] (R) at (3.85,3.3) {$\{abcd\}$};
  \node[tag, anchor=west] at ($(R.east)+(0.18,0)$) {%
    $\sigma:$\\
    $a \rightarrow a$ \\%
    $b \rightarrow b$ \\%
    $c \rightarrow c$ \\%
    $d \rightarrow d$ \\%
    $e \rightarrow \varepsilon$};

  \node[lat-empty] (L1) at (-5.7,3.75) {$\cdots$};
  \node[tag, anchor=east] at ($(L1.west)+(-0.18,0)$) {%
    $\{a,b,c,d,e\} \rightarrow a$};
  \draw[dashed,->,thick] (R.west) --
    node[arrlbl] {$\{aaaa\}$}
    (L1.east);

  \node[lat-empty] (L2) at (-5.7,2.5) {$\cdots$};
  \node[tag, anchor=east] at ($(L2.west)+(-0.18,0)$) {%
    $\{a,b,c,e\} \rightarrow a$ \\%
    $d \rightarrow d$};
  \draw[dashed,->,thick] (R.west) --
    node[arrlbl] {$\{aaa, ad\}$}
    (L2.east);

  \node at (-5.7,1.6) {$\vdots$};

  \node[lat-empty] (D3) at (-5.7,0.92) {$\cdots$};
  \node[tag, anchor=east] at ($(D3.west)+(-0.18,0)$) {%
    $\{a,c,e\} \rightarrow a$ \\%
    $b \rightarrow b$ \\%
    $d \rightarrow d$};
  \draw[dashed,->,thick] (R.west) --
    node[arrlbl] {$\{aba, ad\}$}
    (D3.east);

  \node at (-5.0,0.05) {$\vdots$};

  \node[lat-empty] (D4) at (-4.6,-0.83) {$\cdots$};
  \node[tag, anchor=east] at ($(D4.west)+(-0.18,0)$) {%
    $a \rightarrow a$ \\%
    $\{b,c,e\} \rightarrow b$ \\%
    $d \rightarrow d$};
  \coordinate (RsW) at ($(R.south west)!0.55!(R.west)$);
  \draw[dashed,->,thick] (RsW) --
    node[arrlbl] {$\{abb, bbd\}$}
    (D4.north east);

  \node at (-2.6,-0.83) {$\cdots$};

  \node[lat-empty] (D1) at (-0.5,-0.83) {$\cdots$};
  \node[tag, anchor=west] at ($(D1.east)+(0.18,0)$) {%
    $a \rightarrow a$ \\%
    $\{b, e\} \rightarrow b$ \\%
    $c \rightarrow c$ \\%
    $d \rightarrow d$};
  \draw[dashed,->,thick] (R.south west) --
    node[arrlbl] {$\{ab, bc, d\}$}
    (D1.north east);

  \node[lat-empty] (D2) at (3.85,-0.83) {$\cdots$};
  \node[tag, anchor=west] at ($(D2.east)+(0.18,0)$) {%
    $\{a, e\} \rightarrow a$ \\%
    $b \rightarrow b$ \\%
    $c \rightarrow c$ \\%
    $d \rightarrow d$};
  \draw[dashed,->,thick] (R.south) --
    node[arrlbl] {$\{ab, c, d\}$}
    (D2.north);
\end{tikzpicture}
\end{center}
\label{fig:nonltl}
\caption{Нетривиальный перенос уравнений из не-LTL свойства в несравнимые по морфизму свойства}
\end{figure}

Как было отмечено ранее, ни один из переносов в вышеприведенном примере не может быть совершен между данным и свойством со сравнимым с ним меньшим морфизмом.
Для сравнимых стирающих морфизмов процесс кросс-переносов схож с кросс-переносом для LTL-свойств.

При этом из-за стирания структура переносимого уравнения более непредсказуема и зависит от конкретной системы уравнений (см. рисунок~\ref{fig:cross-ltl-example}).

\begin{figure}[htb]
\begin{center}
\begin{tikzpicture}[
    y=1.38cm,
    lat/.style={
      rectangle,
      draw,
      thick,
      rounded corners=2pt,
      minimum width=2.3cm,
      minimum height=0.95cm,
      align=center,
      inner sep=4pt,
      fill=white
    },
    lat-empty/.style={
      rectangle,
      draw=black!35,
      thick,
      rounded corners=2pt,
      minimum width=1.65cm,
      minimum height=0.72cm,
      align=center,
      inner sep=2pt,
      fill=gray!14,
      text=black!45
    },
    arrlbl/.style={
      fill=white,
      inner sep=2pt,
      font=\footnotesize,
      pos=0.5,
      above
    },
    tag/.style={
      rectangle,
      draw,
      thick,
      rounded corners=1pt,
      font=\scriptsize,
      inner sep=3pt,
      align=center,
      fill=white
    }
  ]
  \node[lat] (D1) at (-2.8,2.2) {$\{abc\}$};
  \node[tag, anchor=east] at ($(D1.west)+(-0.18,0)$) {%
    $\sigma_1:$ \\%
    $a \rightarrow a$ \\%
    $b \rightarrow b$ \\%
    $c \rightarrow c$ \\%
    $\{d, e\} \rightarrow \varepsilon$};

  \node[lat-empty] (D3) at (2.8,2.2) {$\cdots$};
  \node[tag, anchor=west] at ($(D3.east)+(0.18,0)$) {%
    $\sigma_2:$ \\%
    $a \rightarrow a$ \\%
    $\{b, e\} \rightarrow \varepsilon$ \\%
    $\{c, d\} \rightarrow c$};

  \node[lat-empty] (D4) at (0,0) {$\cdots$};
  \node[tag, anchor=east] at ($(D4.west)+(-0.18,0)$) {%
    $a \rightarrow a$ \\%
    $b \rightarrow b$ \\%
    $c \rightarrow c$ \\%
    $d \rightarrow d$ \\%
    $e \rightarrow \varepsilon$};

  \draw[dashed,->,thick] (D1.east) --
    node[arrlbl] {$\{ac\}$}
    (D3.west);
  \draw[dashed,->,thick] (D1.south east) to[bend left=22]
    node[arrlbl] {$\{a,\, b,\, c\}$}
    (D4.north);

  \draw[thick] (D1.south) -- (D4.north);
  \draw[thick] (D3.south) -- (D4.north);
\end{tikzpicture}
\end{center}
\label{fig:cross-ltl-example}
\caption{Кросс-перенос из SP стирающего стандартного морфизма}
\end{figure}

В данном случае в свойство ниже по решетке морфизмов не может сохранить ту же информацию, что и кросс-перенос (даже для двух не-LTL свойств).
В этом их отличие от свойств, задающих LT языки.

Перенос в LTL-свойства из не-LTL-свойств создает трудности для редукции строкового объекта в целом: именно LTL-свойства в некотором
смысле более устойчивы к сохранению информации при строковых операциях, обобщении и т.п. Появление порядка $2^{|A|}$ таких <<спутников>>
создает сложность для анализа в целом, поскольку, если стремиться к точности, необходимо отслеживать также все эти свойства.


\begin{Proposition}
Для LTL-свойств $SP_1 = (S_1, \sigma_1)$ и $SP_2 = (S_2, \sigma_2)$ верно следующее: если уравнение $e \in S_2$ может быть перенесено в $SP_1$, 
то $e$ может быть перенесено в $(S_{meet}, \sigma_3 = \sigma_1 \sqcap \sigma_2)$. 

При этом трансформация строки $w$ уравнения $e$ сохраняется: $\sigma_1(\sigma_2^{-1}(w)) = \sigma_1(\sigma_3(\sigma_2^{-1}(w)))$.
\end{Proposition}

\begin{proof}
    
    Пусть $\sigma_1(\sigma_2^{-1}(w)) = u$ --- слово из уравнения, которое переносится в $SP_1$. 

    $\sigma_3(\sigma_2^{-1}(w)) = \min [ \sigma_1^{-1}(\sigma_1(\sigma_2^{-1}(w))) \cap \sigma_2^{-1}(\sigma_2(\sigma_2^{-1}(w))) ] = 
    \min(\sigma_1^{-1}(u) \cap \sigma_2^{-1}(w))$ и т.к. из посылки теоремы $\sigma_1^{-1}(u) \supseteq \sigma_2^{-1}(w)$, то
    $\sigma_3(\sigma_2^{-1}(w)) = \min(\sigma_2^{-1}(w)) = \sigma_2(w)$.

    Дополнительное утверждение не требует отдельного обоснования, т.к. $\sigma_3 \preceq \sigma_1$ по определению.
\end{proof}

% TODO: раздел об обработке строкового объекта с различным набором морфизмов

\anonsection{СПИСОК ИСПОЛЬЗОВАННЫХ ИСТОЧНИКОВ}
\printbibliography[heading=none]


\end{document}